\documentclass[11pt,a4paper]{article}

\usepackage[utf8]{inputenc}
\usepackage[T1]{fontenc}
\usepackage{lmodern}
\usepackage{geometry}
\usepackage{amsmath,amssymb}
\usepackage{graphicx}
\usepackage{booktabs}
\usepackage{hyperref}
\usepackage{enumitem}
\usepackage{setspace}

\geometry{margin=2.5cm}
\onehalfspacing

\title{
Link Prediction in Citation Networks using Graph Neural Networks\\
\large Master-Level Graph Neural Networks Course Project Proposal
}

\author{
Sándor Baranyi \\ \texttt{ct9xfj@inf.elte.hu}
\and
Balázs Bordács \\ \texttt{adjdme@inf.elte.hu}
\and
Máté Kovács \\ \texttt{u5bky4@inf.elte.hu}
\and
Bence Adorján Nagy-Mohos \\ \texttt{d5vd5e@inf.elte.hu}
}

\date{\today}

\begin{document}

\maketitle

\begin{abstract}

This project investigates the use of Graph Neural Networks (GNNs) for link prediction
in citation networks using benchmark datasets such as Cora and PubMed.
The goal is to model citation graphs and recommend relevant papers to cite for a given paper.
We compare classical link prediction methods with neural approaches,
including Graph Convolutional Networks and GraphSAGE.
Additional experiments investigate the effect of structural encodings and
directed versus undirected graph modeling.
The project focuses on reproducible, CPU-friendly experiments using standard dataset splits.

\end{abstract}


\section{Introduction}

Graph-structured data appears in many real-world applications,
including social networks, biological networks, and citation networks.
Citation graphs represent relationships between scientific papers,
where nodes correspond to papers and edges represent citation links.

Predicting missing links in citation networks can be interpreted as
recommending relevant references for a given paper.
This task is important for recommendation systems, knowledge graph completion,
and scientific information retrieval.

Graph Neural Networks (GNNs) have recently shown strong performance
on graph learning tasks.
In this project, we investigate whether simple GNN architectures
can outperform classical link prediction methods
on citation graph benchmarks.


\section{Problem Definition}

We model a citation network as a graph

\[
G = (V, E)
\]

where

\begin{itemize}
\item $V$ is the set of papers
\item $E$ is the set of citation edges
\end{itemize}

The task is link prediction.

Given two nodes $u$ and $v$, the goal is to predict
whether a citation edge should exist between them.

This can be formulated as a binary classification problem
over node pairs.


\section{Objectives and Research Questions}

The main objective of the project is to evaluate the effectiveness
of Graph Neural Networks for link prediction in citation networks.

Research questions:

\begin{itemize}
\item Do GNN-based methods outperform classical link prediction methods?
\item Does structural encoding improve link prediction accuracy?
\item Does modeling the graph as directed improve performance?
\item Can simple GNN models achieve good results without GPU training?
\end{itemize}


\section{Novelty and Motivation}

Citation datasets are commonly used for node classification,
but link prediction is less explored in course-level projects.

This project focuses on:

\begin{itemize}
\item link prediction instead of node classification
\item comparison of classical and neural methods
\item structural and positional encodings
\item directed vs undirected citation graphs
\end{itemize}

The goal is to analyze how much benefit GNNs provide
over simpler graph algorithms.


\section{Dataset Audit}

We use standard benchmark citation datasets.

\subsection{Cora}

\begin{itemize}
\item Nodes: 2708
\item Edges: 5429
\item Features: bag-of-words
\item Classes: 7
\end{itemize}

\subsection{PubMed}

\begin{itemize}
\item Nodes: 19717
\item Edges: 44338
\item Features: TF-IDF vectors
\end{itemize}

Datasets are loaded using PyTorch Geometric
with standard train/test splits.

Limitations:

\begin{itemize}
\item small graph size
\item sparse edges
\item limited feature information
\end{itemize}


\section{Literature Review}

The project builds on the following works.

\begin{itemize}

\item Kipf and Welling — Graph Convolutional Networks

\item Hamilton et al. — GraphSAGE

\item Grover and Leskovec — Node2Vec

\item Page et al. — PageRank

\item Zhang and Chen — Link Prediction with Graph Neural Networks

\item Dwivedi et al. — Benchmarking Graph Neural Networks

\end{itemize}

These papers provide baseline methods
and motivation for the proposed experiments.


\section{Baseline Methods}

We implement classical link prediction methods.

\begin{itemize}
\item Common Neighbors
\item Adamic-Adar
\item PageRank
\item Node2Vec embeddings
\end{itemize}

These serve as non-neural baselines.


\section{GNN Models}

We implement simple Graph Neural Networks.

\begin{itemize}
\item Graph Convolutional Network (GCN)
\item GraphSAGE
\end{itemize}

Implementation will use PyTorch Geometric.


\section{Research Extensions}

We plan to experiment with:

\begin{itemize}
\item positional encodings
\item structural features
\item directed vs undirected graphs
\item different embedding sizes
\item different number of layers
\end{itemize}


\section{Evaluation}

Metrics:

\begin{itemize}
\item Accuracy
\item AUC
\item Precision
\item Recall
\item F1 score
\end{itemize}

We use fixed dataset splits
to ensure reproducibility.


\section{Reproducibility Plan}

\begin{itemize}
\item GitHub repository
\item fixed random seeds
\item CPU-friendly experiments
\item PyTorch Geometric
\item documented environment
\end{itemize}


\section{Work Plan}

\begin{tabular}{ll}
Week 5 & Proposal \\
Week 6 & Dataset audit \\
Week 7 & Baselines \\
Week 8 & GNN implementation \\
Week 9 & Experiments \\
Week 10 & Final report \\
\end{tabular}


\section{Risks and Limitations}

\begin{itemize}
\item small datasets
\item overfitting
\item implementation complexity
\item time limitations
\end{itemize}


\section{Conclusion}

This project investigates link prediction in citation networks
using classical graph algorithms and Graph Neural Networks.
The goal is to evaluate the usefulness of GNNs
in a reproducible and CPU-friendly experimental setup.


\begin{thebibliography}{9}

\bibitem{gcn}
Kipf, T. N., Welling, M.
Semi-Supervised Classification with Graph Convolutional Networks.

\bibitem{sage}
Hamilton, W., Ying, Z., Leskovec, J.
Inductive Representation Learning on Large Graphs.

\bibitem{node2vec}
Grover, A., Leskovec, J.
node2vec: Scalable Feature Learning for Networks.

\bibitem{pagerank}
Page, L. et al.
The PageRank Citation Ranking.

\bibitem{seal}
Zhang, M., Chen, Y.
Link Prediction Based on Graph Neural Networks.

\bibitem{benchmark}
Dwivedi, V. et al.
Benchmarking Graph Neural Networks.

\end{thebibliography}

\end{document}